\songtitle{Der Quacksalber}{Schandmaul}{2022}

\adjustboxDefault{%
\guitarchord{G}
\guitarchord{D}
\guitarchord{A}
\guitarchord{Bm}}
\strummingTwoBeatsAndEightMovementsPerBeat{d.du.udud.du.udu}

\begin{guitar}
	\songsection{Intro}
	{\footnotesize\begin{tabular}{|l|l|l|l|l}
			G & D & A & Bm & (x5) \\
			G & D & A & A & 
	\end{tabular}}
	
	\songsection{Strophe 1}
	Packt dich die [G]Angst,[D] dass man [A]dir ans Leder [Bm]will?
	Ein Mörder [G]frech mischt hi[D]nein ein [A]Gift dir in den [Bm]Wein?
	Hab' ich [G]hier aus fernen [D]Landen ein gar [A]wunderliches [Bm]Horn.
	Tust [G]du daraus dich [D]gütlich, bleibst du [A]ungeschor'n.
	
	Oder vergisst du allzu oft sehr viel?
	Trag diesen Schmuck um den Hals. Er hilft dir jedenfalls.
	Getrocknet ist die Zunge, die war einst dem Wiedehopf.
	An der Kette ziert sie dir vortrefflich deinen Schopf.
	
	\begin{highlightbar}
		\songsection{Refrain}
		[G]Salben, Tränke [D]und Tinkt[A]uren
		Aus [G]Dingen, die wir [D]nie erf[A]uhren,
		Und [G]eine Prise [D]Zaubere[Bm]i. Das war die [A]Quacksalber[G]ei.[D A]{}
	\end{highlightbar}
	
	\songsection{Instrumental}
	{\footnotesize\begin{tabular}{|l|l|l|l|l}
			G & D & A & Bm & (x3) \\
			G & D & A & A & 
	\end{tabular}}
	
	\songsection{Strophe 2}
	Nimm diesen Saft aus Tausendgüldenkraut
	Und noch den Wermut dazu. Vermische sie im Nu.
	Und dann koch dir daraus Tee, den du genießt um Mitternacht.
	Denkst du dabei an Jungfrau'n, er den Darm dir leichter macht.
	
	Die Mannespracht dir schon lang nicht mehr erwacht,
	Dann bist du bei mir am Ziel. Schmier vor dem Liebesspiel
	Brennnessel- und Honigsalbe direkt auf den Schwanz.
	Die Weiber werden Schlange stehen und sagen: "Ja, der kann's."
	
	\begin{highlightbar}
		\songsection{Refrain} \songsnippet{Salben, Tränke und Tinkturen...} \optionalChord{(x2)}
	\end{highlightbar}
	
	\songsection{Instrumental}
	{\footnotesize\begin{tabular}{|l|l|l|l|l}
			G & D & A & Bm & (x7) \\
			G & D & A & A & D
	\end{tabular}}
\end{guitar}